%% ============================================================
%%  paper3_kappa_rate.tex
%%  "The κ-Rate: A Riba-Free Monetary Alternative
%%   Derived from Perpetual Contract Theory"
%%  Ahmed, Bhuyan & Islam — February 2026
%% ============================================================
\documentclass[10pt,twocolumn]{article}
\usepackage[utf8]{inputenc}
\usepackage[margin=0.75in]{geometry}
\usepackage{amsmath,amssymb,amsthm}
\usepackage{graphicx}
\usepackage{hyperref}
\usepackage[numbers]{natbib}
\usepackage{booktabs}
\usepackage{abstract}
\usepackage{xurl}
\usepackage{titlesec}
\usepackage{authblk}
\usepackage{xcolor}
\usepackage{array}
\usepackage[protrusion=true,expansion=false]{microtype}
\usepackage{enumitem}
\usepackage{caption}
\usepackage{multirow}
\usepackage{bm}

%% Section spacing
\titlespacing*{\section}{0pt}{13pt}{6pt}
\titlespacing*{\subsection}{0pt}{10pt}{4pt}
\titlespacing*{\subsubsection}{0pt}{8pt}{3pt}

%% Abstract style
\renewcommand{\abstractnamefont}{\normalfont\bfseries\large}
\renewcommand{\abstracttextfont}{\normalfont\small}

%% Theorem environments
\theoremstyle{plain}
\newtheorem{theorem}{Theorem}
\newtheorem{proposition}{Proposition}
\newtheorem{corollary}{Corollary}
\newtheorem{lemma}{Lemma}
\theoremstyle{definition}
\newtheorem{definition}{Definition}
\newtheorem{remark}{Remark}
\newtheorem{example}{Example}

%% Math operators
\DeclareMathOperator{\E}{\mathbb{E}}
\DeclareMathOperator*{\argmin}{arg\,min}
\newcommand{\Qa}{Q^a}
\newcommand{\R}{\mathbb{R}}
\newcommand{\Prob}{\mathbb{P}}
\newcommand{\F}{\mathcal{F}}
\newcommand{\krate}{\hat{\kappa}}

\definecolor{ribared}{RGB}{180,30,30}
\definecolor{halalgreen}{RGB}{30,120,60}
\definecolor{noteblue}{RGB}{20,60,160}

\hypersetup{colorlinks=true,linkcolor=blue,citecolor=blue,urlcolor=blue}

%% ============================================================
\title{\LARGE\textbf{The $\kappa$-Rate:\\
A Riba-Free Monetary Alternative\\
Derived from Perpetual Contract Theory}}

\author[1]{Shehzad Ahmed}
\author[2]{Rafiqul Bhuyan}
\author[3]{Rubaiyat Islam}

\affil[1]{Department of Finance, Independent University, Bangladesh\\
\small\texttt{2120922@iub.edu.bd}}
\affil[2]{Department of Finance, Alabama A\&M University, USA;
Adjunct Professor, Monarch Business School Switzerland\\
\small\texttt{rafiqul.bhuyan@aamu.edu}}
\affil[3]{Department of Software Engineering, Daffodil International University,
Bangladesh\\ \small\texttt{islam.swe@diu.edu.bd}}

\date{February 2026}
%% ============================================================

\begin{document}

\twocolumn[\begin{@twocolumnfalse}
\maketitle\thispagestyle{empty}

\begin{abstract}
\noindent
We propose the \textbf{$\kappa$-rate} --- the convergence intensity of
perpetual futures contracts \cite{Ackerer2024} --- as the first
mathematically rigorous, no-arbitrage-grounded, riba-free alternative
to the risk-free interest rate $r$ in Islamic monetary economics.

The $\kappa$-rate is: (i)~\emph{endogenous} --- market-implied from
the observed perpetuals basis in continuous time; (ii)~\emph{observable}
--- computable from two market prices and publishable on-chain in real
time; (iii)~\emph{riba-free} --- an event intensity, not a reward for
the time value of money; and (iv)~\emph{complete} --- sufficient to
construct a full term structure (the $\kappa$-yield curve) analogous
to the conventional yield curve.

We develop the theoretical foundations in four steps. First, we
decompose the classical interest rate $r$ (Fisher \cite{Fisher1930},
Wicksell \cite{Wicksell1898}, Keynes \cite{Keynes1936}) into its
components and identify the riba component as the pure time-preference
premium, which the Ackerer-Hugonnier-Jermann no-arbitrage theorem
proves is eliminable ($\iota = 0$). Second, we show that the $\kappa$-rate
is the Wicksellian ``natural rate'' (a rate grounded in real economic
dynamics rather than monetary policy) made rigorous and riba-free for
the first time. Third, we construct the $\kappa$-yield curve
$\kappa(T) = 1/T$ as the Islamic analog of the CIR
\cite{CIR1985} term structure, grounded in exponential stopping time
distributions. Fourth, we demonstrate that $\kappa$ can serve all
monetary functions currently performed by $r$ in the Islamic context:
pricing credit instruments (Section~\ref{sec:credit}), anchoring
actuarial takaful contributions (Section~\ref{sec:takaful}), guiding
Islamic monetary policy (Section~\ref{sec:policy}), and providing a
sovereign benchmark rate for sukuk issuance (Section~\ref{sec:sovereign}).

Empirical support comes from the Baraka Protocol --- a Shariah-compliant
perpetual futures exchange deployed on Arbitrum Sepolia (chainId 421614)
--- where $\kappa$ is on-chain observable in real time. The contemporary
Shariah ruling of SeekersGuidance \cite{SeekersGuidance2025} validates
that $\iota = 0$ resolves the riba objection, providing jurisprudential
grounding for the $\kappa$-rate as a Shariah-compatible monetary instrument.

\medskip
\noindent\textbf{Keywords:} $\kappa$-rate, riba-free monetary policy,
Islamic monetary economics, natural rate, perpetual contracts, yield
curve, Wicksell, Fisher decomposition, no-arbitrage, $\iota = 0$,
Baraka Protocol
\end{abstract}
\bigskip
\end{@twocolumnfalse}]

%% ============================================================
\section{Introduction}
\label{sec:intro}
%% ============================================================

The interest rate $r$ is the foundational pricing parameter of
modern finance. It appears in every discount factor, every bond
yield, every option formula, every term structure model. It is the
rate at which future cash flows are brought to present value.
Islamic finance prohibits it.

This is not a peripheral prohibition. The Quranic injunction
against \textit{riba} --- ``Allah has permitted trade and
forbidden \textit{riba}'' (Al-Baqarah 2:275) --- is absolute,
applying to all forms of predetermined excess in monetary exchange.
The risk-free rate $r$ is, under any classical decomposition
(Fisher \cite{Fisher1930}, Böhm-Bawerk \cite{BohmBawerk1890},
Keynes \cite{Keynes1936}), precisely the price of time preference
in monetary form: the lender receives $r$ per unit time merely for
parting with money, independent of any real economic outcome. This
is \textit{riba al-nasi'ah} --- the prohibited excess for deferment
--- by all four Sunni schools and the Ja'fari school
\cite{Usmani2002,ElGamal2006}.

The consequence for Islamic finance is catastrophic. Without $r$,
there is no discount factor. Without a discount factor, bonds
cannot be priced. Without bond pricing, there is no credit market.
Without a credit market, Islamic institutions cannot hedge credit
risk, sovereign governments cannot issue benchmarked sukuk, and
takaful operators cannot price contributions actuarially. The
entire credit infrastructure of modern finance --- a market of
roughly \$130 trillion in bonds and \$10 trillion in credit
derivatives \cite{BIS2024} --- is inaccessible to Islamic
participants on Shariah-compliant terms.

Islamic monetary economists have recognised this problem for
five decades (Chapra \cite{Chapra1985}, Khan \& Mirakhor
\cite{KhanMirakhor1987}, Iqbal \& Mirakhor \cite{IqbalMirakhor2011}).
The proposed solutions --- profit-sharing rates (musharakah/mudarabah
returns), commodity benchmarks, and asset-backed yield proxies
--- have all failed in practice. Iqbal \& Molyneux
\cite{IqbalMolyneux2005} document that even in the most
advanced Islamic finance jurisdictions (Malaysia, Bahrain,
UAE), rates charged on Islamic products track LIBOR/SOFR
closely, with the interest content merely re-labelled.
Khan \& Abdallah \cite{KhanAbdallah2017} term this the
``ghost of LIBOR'' --- riba is present in Islamic pricing
even when explicitly prohibited.

The failure has a structural cause: \emph{no prior work has
derived a complete, mathematically rigorous, no-arbitrage-grounded
alternative to $r$ from first principles}.

This paper provides that alternative. We call it the
\textbf{$\kappa$-rate}: the convergence intensity of a perpetual
futures contract, proved by Ackerer, Hugonnier \& Jermann
\cite{Ackerer2024} to achieve unique no-arbitrage pricing at
$\iota = 0$ (zero interest parameter). The $\kappa$-rate is not
a label change. It is a mathematically distinct pricing parameter
with provably no riba content, capable of performing all the
monetary functions currently assigned to $r$ in an Islamic economy.

\subsection*{The Gap in One Sentence}

Wicksell \cite{Wicksell1898} proposed in 1898 that the ``natural
rate'' of interest should be grounded in real economic productivity
rather than monetary convention --- but it remained an interest rate
with time-preference content. Chapra \cite{Chapra1985} proposed
in 1985 that Islamic monetary policy needs a real-activity-based
pricing signal --- but never provided a mathematical formula.
The $\kappa$-rate is the Wicksellian natural rate made rigorous,
riba-free, and on-chain observable --- 128 years after the question
was first posed.

\subsection*{Structure}

Section~\ref{sec:decomposition} decomposes $r$ into components and
identifies the riba component.
Section~\ref{sec:kappa_def} defines the $\kappa$-rate from first
principles.
Section~\ref{sec:kappa_vs_r} contrasts $\kappa$ with $r$ across
all standard properties.
Section~\ref{sec:yieldcurve} constructs the $\kappa$-yield curve.
Section~\ref{sec:applications} develops applications: credit
instruments, takaful, monetary policy, and sovereign benchmarking.
Section~\ref{sec:empirical} presents on-chain empirical evidence
from the Baraka Protocol.
Section~\ref{sec:shariah} provides the Shariah analysis.
Section~\ref{sec:discussion} compares $\kappa$ with prior proposals
and discusses limitations.

%% ============================================================
\section{Decomposing $r$: What Is the Riba Component?}
\label{sec:decomposition}
%% ============================================================

\subsection{The Fisher Decomposition}

Fisher \cite{Fisher1930} established the modern theory of interest
as an equilibrium between two forces:
\begin{enumerate}[leftmargin=*,noitemsep]
\item \textbf{Time preference (impatience):} the subjective
preference of individuals for present over future consumption,
independent of any investment opportunity.
\item \textbf{Investment opportunity:} the marginal productivity
of capital --- the real return from deploying resources in
production.
\end{enumerate}

The equilibrium interest rate $r$ equates marginal impatience
to marginal productivity. Fisher's formulation also yields the
famous Fisher equation:
\begin{equation}
r_{\text{nominal}} = r_{\text{real}} + \pi
\label{eq:fisher}
\end{equation}
where $\pi$ is expected inflation. The market rate therefore
contains three components: real productivity, time preference,
and inflation compensation.

\subsection{Böhm-Bawerk's Distinction}

Böhm-Bawerk \cite{BohmBawerk1890} distinguished two grounds
for interest:
\begin{itemize}[leftmargin=*, noitemsep]
\item The \textbf{productivity ground}: present capital generates
future output exceeding the present input (roundabout production).
This is the return to real productive activity.
\item The \textbf{time preference ground}: present goods are
valued more than future goods by an ``agio'' (premium) for
their presentness --- independent of productivity.
\end{itemize}

Islamic economics accepts the productivity ground (profit from
real activity is halal) but rejects the time preference ground
(a predetermined premium for the passage of time, payable on
a money loan regardless of economic outcome, is riba). This
is the precise jurisprudential boundary.

\subsection{The Riba Component}

Under the Fisher-Böhm-Bawerk decomposition, the risk-free
interest rate $r$ can be written:
\begin{equation}
r = \underbrace{r_{\text{prod}}}_{\text{halal}} +
\underbrace{r_{\text{pref}}}_{\text{riba}} +
\underbrace{\pi}_{\text{neutral}}
\label{eq:riba_decomp}
\end{equation}

The riba component $r_{\text{pref}}$ is the pure time preference
premium: the predetermined, fixed return on a money loan
attributable solely to the passage of time. It is independent of:
\begin{itemize}[leftmargin=*,noitemsep]
\item the real economic outcome of the borrower,
\item the productive use of the capital,
\item market conditions at the time of repayment.
\end{itemize}

Under El-Gamal's \cite{ElGamal2006} operationalization of riba
as ``predetermined excess in exchange, independent of economic
outcome,'' $r_{\text{pref}}$ is riba by definition.

\subsection{The Wicksell Natural Rate and Its Limit}

Wicksell \cite{Wicksell1898} proposed the ``natural rate'' of
interest $r^*$ --- the rate at which investment equals saving
in a real economy, determined by the marginal productivity of
capital without monetary distortion:
\begin{equation}
r^* = \text{MPK} = \frac{\partial F(K,L)}{\partial K}
\label{eq:wicksell}
\end{equation}

This is closer to the Islamic ideal: a rate grounded in real
activity rather than monetary policy. But Wicksell's natural
rate still includes time preference: it is still an interest
rate that compensates capital for time elapsed, just one
calibrated to real productivity. The riba component persists.

\begin{remark}
The $\kappa$-rate achieves what Wicksell attempted but could
not complete. It eliminates $r_{\text{pref}}$ entirely from
the pricing formula, leaving only the productivity-equivalent
component in the form of convergence dynamics.
The $\kappa$-rate is Wicksell's natural rate with
the riba subtracted.
\end{remark}

%% ============================================================
\section{The $\kappa$-Rate: Definition and Properties}
\label{sec:kappa_def}
%% ============================================================

\subsection{From Perpetual Futures to Monetary Rate}

We recall the foundational result. Ackerer, Hugonnier \& Jermann
\cite{Ackerer2024} prove that the unique no-arbitrage price of a
perpetual futures contract on a spot asset $X_t$ (under constant
parameters) is:
\begin{equation}
f_t = \frac{\kappa}{\kappa - \iota + r_b - r_a} \cdot X_t
\label{eq:perp_price}
\end{equation}
where $\kappa > 0$ is the \textbf{convergence intensity}, $\iota
\geq 0$ is the interest parameter, and $r_a, r_b$ are the
long and short funding rates. Theorem~3 of Ackerer et al.\
proves that $\iota = 0$ satisfies no-arbitrage. At $\iota = 0$
and $r_a = r_b$:
\begin{equation}
f_t = X_t
\label{eq:parity}
\end{equation}
The perpetual futures price equals the spot price identically
--- with the convergence dynamics entirely captured by $\kappa$.

The $\kappa$-rate is implied from observed market prices:
\begin{equation}
\boxed{\hat\kappa_t = \frac{f_t/X_t \cdot (r_b - r_a)}{1 - f_t/X_t}}
\label{eq:kappa_implied}
\end{equation}
This is the \textbf{market-implied $\kappa$-rate}: a continuous,
on-chain observable derived from two market prices, with no
interest rate as an input.

\subsection{Formal Definition}

\begin{definition}[$\kappa$-Rate]
\label{def:kappa}
The \textbf{$\kappa$-rate} for an instrument linking a financial
claim to an underlying real asset $X_t$ is the convergence
intensity $\kappa > 0$ that:
\begin{enumerate}[leftmargin=*,noitemsep]
\item Parameterizes the distribution of the stopping time $\theta_t
\sim t + \mathrm{Exp}(\kappa)$ at which the claim is settled.
\item Is implied endogenously from market prices via
equation~\eqref{eq:kappa_implied}.
\item Satisfies $\iota = 0$: no interest parameter is charged.
\item Achieves unique no-arbitrage pricing by Theorem~3 of
Ackerer et al.\ \cite{Ackerer2024}.
\end{enumerate}
\end{definition}

\subsection{Economic Interpretation of $\kappa$}

$\kappa$ measures the \textbf{speed of convergence} between a
financial instrument's price and its underlying real asset value.
Specifically:
\begin{itemize}[leftmargin=*,noitemsep]
\item \textbf{High $\kappa$:} Rapid mean-reversion. The perpetual
price tracks spot closely. Short expected stopping time
$\E[\theta_t - t] = 1/\kappa$. The market is confident about
near-term convergence.
\item \textbf{Low $\kappa$:} Slow convergence. Wide basis.
Longer expected stopping time. The market perceives structural
uncertainty about when the real asset will reach its fair value.
\item \textbf{$\kappa \to \infty$:} Instantaneous convergence;
$f_t = X_t$ always. Perfect tracking.
\item \textbf{$\kappa \to 0$:} No convergence; the perpetual
becomes untethered from the underlying. The instrument loses
economic meaning.
\end{itemize}

This interpretation is fundamentally different from $r$:
\begin{center}
\small
\begin{tabular}{@{}lll@{}}
\toprule
\textbf{Property} & \textbf{$r$ (interest rate)} & \textbf{$\kappa$ ($\kappa$-rate)} \\
\midrule
Determined by & Central bank & Market prices \\
Measures & Time preference & Convergence speed \\
Riba content & Yes ($r_{\text{pref}} > 0$) & No \\
Real basis & Partial & Complete \\
Negative values & Possible ($r < 0$) & Impossible ($\kappa > 0$) \\
Observable & Policy setting & On-chain, real-time \\
\bottomrule
\end{tabular}
\end{center}

%% ============================================================
\section{$\kappa$ vs.\ $r$: A Complete Comparison}
\label{sec:kappa_vs_r}
%% ============================================================

\subsection{The Riba Test}

El-Gamal's \cite{ElGamal2006} criterion for riba: a payment
is riba if it is (1)~predetermined, (2)~in excess of principal,
and (3)~independent of real economic outcome.

\begin{proposition}[The $\kappa$-Rate is Not Riba]
\label{prop:not_riba}
The $\kappa$-rate satisfies none of the three riba criteria:
\begin{enumerate}[leftmargin=*,noitemsep]
\item \textbf{Not predetermined.} $\kappa_t$ fluctuates
continuously with market conditions. It is implied from
market prices, not contractually fixed.
\item \textbf{Not excess of principal.} The $\kappa$-rate
does not generate a return over principal. At $\iota = 0$
and $r_a = r_b$, the instrument price equals the underlying
asset value ($f_t = X_t$) --- no excess.
\item \textbf{Not independent of economic outcome.}
$\kappa$ is backed out from the observed basis $f_t / X_t$,
which reflects real market views about asset convergence. If
the underlying asset performs well, the basis narrows (high
$\kappa$); if poorly, the basis widens (low $\kappa$). The
$\kappa$-rate is tied to real economic dynamics.
\end{enumerate}
\end{proposition}

\begin{proof}
(1) follows from equation~\eqref{eq:kappa_implied}: $\hat\kappa_t$
is a continuous function of market prices $f_t, X_t$, both of
which vary. (2) follows from equation~\eqref{eq:parity}: at
$\iota = 0$ and $r_a = r_b$, $f_t = X_t$. (3) follows from the
market-determination of $\hat\kappa_t$: it is implied from
observed prices, not set exogenously.\qquad$\square$
\end{proof}

\subsection{Replacing $r$ in Standard Pricing Formulas}

The key question is whether $\kappa$ can perform the mathematical
functions of $r$. We demonstrate this across three standard
settings.

\textbf{Bond pricing.}
A conventional zero-coupon bond: $P = e^{-rT}$.
The $\kappa$-equivalent (Islamic perpetual sukuk with expected
maturity $T = 1/\kappa$):
\begin{equation}
P^{\kappa} = \E^Q[X_\tau] = X_t, \quad \tau\sim\mathrm{Exp}(\kappa)
\label{eq:kappa_bond}
\end{equation}
The discount comes not from $e^{-rT}$ but from the stopping time
distribution: the probability that the redemption event is more
than $T$ periods away decays exponentially at rate $\kappa$,
achieving the same economic effect without any interest.

\textbf{Option pricing.}
In Black-Scholes \cite{BlackScholes1973}, the risk-free rate
$r$ appears as the discount factor and in the drift of the
risk-neutral measure. At $\iota = 0$ and $r_a = r_b$, the
everlasting put from Ackerer et al.\ Proposition~6:
\begin{equation}
\Pi(x,K) = \frac{K^{1-\beta_p}}{1-\beta_p}\cdot x^{\beta_p},
\quad \beta_p = \tfrac{1}{2}
- \sqrt{\tfrac{1}{4} + \tfrac{2\kappa}{\sigma^2}}
\label{eq:evput}
\end{equation}
replaces $r$ with $\kappa$ entirely in the exponent.

\textbf{Credit default swap.}
A conventional CDS spread $s$ solves $s = r + \lambda - r_{\text{risk-free}}
\approx \lambda$ when the risk-free rate is small.
The iCDS spread \cite{Ahmed2026b}: $\kappa_{\text{iCDS}} = \lambda$
(the credit hazard rate), with $\iota = 0$, achieves the same
risk-transfer without the interest baseline.

\subsection{Policy Functions of $r$ and Their $\kappa$ Analogs}

Table~\ref{tab:functions} lists the monetary policy functions of
$r$ and the corresponding $\kappa$-rate mechanism.

\begin{table}[h!]
\centering
\caption{Monetary functions: $r$ vs.\ $\kappa$.}
\label{tab:functions}
\small
\begin{tabular}{@{}p{2.8cm}p{2.1cm}p{2.1cm}@{}}
\toprule
\textbf{Function} & \textbf{$r$ mechanism} & \textbf{$\kappa$ mechanism} \\
\midrule
Price of credit risk & $r + \lambda$ spread & $\kappa \equiv \lambda$ at $\iota=0$ \\
Discount factor & $e^{-rT}$ & $\mathrm{Exp}(\kappa)$ survival prob. \\
Term structure & Yield curve $r(T)$ & $\kappa(T) = 1/T$ \\
Monetary signal & Central bank sets $r$ & Market implies $\hat\kappa$ \\
Sovereign benchmark & Government bond yield & Perpetual sukuk implied $\kappa$ \\
Insurance pricing & $e^{-r} \cdot \text{claim}$ & $\Pi(x,K) = C_p x^{\beta_p}$ \\
\bottomrule
\end{tabular}
\end{table}

%% ============================================================
\section{The $\kappa$-Yield Curve}
\label{sec:yieldcurve}
%% ============================================================

\subsection{Construction}

The conventional yield curve $r(T)$ maps maturity $T$ to the
yield on a zero-coupon risk-free bond. It is the term structure
of interest rates, grounded in the expectation hypothesis and
risk premia over time. Cox, Ingersoll \& Ross \cite{CIR1985}
provide the canonical equilibrium derivation; Vasicek
\cite{Vasicek1977} provides the Gaussian version; Nelson \&
Siegel \cite{NelsonSiegel1987} provide the empirical estimation
framework.

The \textbf{$\kappa$-yield curve} is the Islamic analog: the
mapping $\kappa(T)$ of convergence intensities from perpetual
instruments of different expected horizons.

\begin{proposition}[$\kappa$-Yield Curve]
\label{prop:kappa_curve}
For an instrument with expected horizon $T$ (the expected time
until the stopping event under no other information), the
$\kappa$-yield curve satisfies:
\begin{equation}
\kappa(T) = \frac{1}{T}, \qquad T > 0
\label{eq:kappa_curve}
\end{equation}
The curve is monotonically decreasing in $T$: short-horizon
instruments have high $\kappa$ (rapid convergence), long-horizon
instruments have low $\kappa$ (slow convergence).
\end{proposition}

\begin{proof}
For a stopping time $\theta_t \sim \mathrm{Exp}(\kappa)$,
the expected stopping time is $\E[\theta_t - t] = 1/\kappa$.
Setting this equal to the instrument's target horizon $T$
gives $\kappa = 1/T$.\qquad$\square$
\end{proof}

\subsection{Comparison with the Conventional Yield Curve}

Table~\ref{tab:yield_curve} contrasts the two term structures.

\begin{table}[h!]
\centering
\caption{$\kappa$-yield curve vs.\ conventional yield curve.}
\label{tab:yield_curve}
\small
\begin{tabular}{@{}p{2.6cm}p{2.0cm}p{2.4cm}@{}}
\toprule
\textbf{Property} & \textbf{CIR/Vasicek} & \textbf{$\kappa$-curve} \\
\midrule
Pricing parameter & Short rate $r_t$ & Conv.\ intensity $\kappa_t$ \\
Riba content & Yes & No \\
Central bank role & Sets $r$ & None (market-implied) \\
Shape & Any & Monotone decreasing \\
Calibration & Rate expectations & Asset timelines \\
Negative rates & Possible & Impossible \\
DeFi observable & No & Yes (on-chain) \\
\bottomrule
\end{tabular}
\end{table}

\subsection{Empirical Estimation}

By Proposition~\ref{prop:kappa_curve}, to estimate the
$\kappa$-yield curve one needs:
\begin{enumerate}[leftmargin=*,noitemsep]
\item Observed perpetuals basis $f_t^{(i)}/X_t^{(i)}$ for
instruments $i = 1,\ldots,n$ with different underlying asset
horizons.
\item Implied $\hat\kappa_t^{(i)}$ from equation~\eqref{eq:kappa_implied}.
\item The Nelson-Siegel \cite{NelsonSiegel1987} functional form
applied to the $(\hat\kappa^{(i)}, T^{(i)})$ pairs to produce
a smooth curve.
\end{enumerate}

For the Baraka Protocol (Section~\ref{sec:empirical}), the
single BTC perpetual instrument provides the spot
$\hat\kappa_t$; a term structure requires perpetuals on
multiple RWA underlyings with different expected lifetimes
(commodity contracts, infrastructure sukuk, mortgage-equivalent
instruments).

%% ============================================================
\section{Applications of the $\kappa$-Rate}
\label{sec:applications}
%% ============================================================

\subsection{Islamic Credit Instruments}
\label{sec:credit}

The $\kappa$-rate provides the pricing foundation for a complete
riba-free credit market. The formal proof (the isomorphism between
the Ackerer framework and the Duffie-Singleton credit model) is
established in Ahmed, Bhuyan \& Islam \cite{Ahmed2026b}. We
summarize the pricing formulas.

\textbf{Perpetual sukuk.} Backed by a real asset with value $X_t$
and expected maturity $T$:
\begin{equation}
P_t^{\text{sukuk}} = \E^Q[X_\tau] = X_t
\quad (r_a = r_b, \iota = 0)
\label{eq:sukuk}
\end{equation}
The $\kappa$-rate calibration: $\kappa = 1/T$ from the asset's
expected project timeline (engineering estimates, actuarial tables,
regulatory schedules).

\textbf{iCDS (Islamic Credit Default Swap).}
Protection buyer pays $\kappa\,dt$ continuously; seller pays
$X_\tau$ at credit event $\tau$:
\begin{equation}
V_t^{\text{iCDS}} = \E^Q[X_\tau] - \kappa\int_t^\infty
e^{-\kappa(s-t)}\,ds = X_t - 1
\label{eq:icds}
\end{equation}
for normalized sukuk price (par = 1). The iCDS $\kappa$ is
calibrated from the reference entity's historical default
frequency and rating transition matrices.

\textbf{Sovereign perpetual sukuk.}
A government issues sukuk backed by an infrastructure project
(expected completion $T$ years, value $X_t$). Secondary market
price: $P_t = X_t$ (at-par, riba-free). The government pays
continuous funding $\pi_t = \kappa(f_t - X_t)\,dt$ to holders.
No SOFR benchmark required.

\subsection{Takaful (Islamic Insurance)}
\label{sec:takaful}

Billah \cite{Billah2007} identifies the core takaful pricing
tension: actuarial fairness requires discounting expected claims,
but standard discounting uses $r$ (riba). The $\kappa$-rate
resolves this.

The actuarially fair takaful contribution for floor protection
at $K$ on asset $X_t$ (volatility $\sigma$, claim intensity $\kappa$):
\begin{equation}
\Pi_{\text{takaful}}(x,K) = \frac{K^{1-\beta_p}}{1-\beta_p}
\cdot x^{\beta_p},\quad
\beta_p = \frac{1}{2} - \sqrt{\frac{1}{4} + \frac{2\kappa}{\sigma^2}}
\label{eq:takaful}
\end{equation}
No interest rate appears. The $\kappa$-rate calibration uses
actuarial tables (historical claim frequency $=\kappa$). The
contribution~\eqref{eq:takaful} is provably no-arbitrage
\cite{Ahmed2026b}, resolving the Htay-Salman circularity
\cite{HtaySalman2012}.

\subsection{Islamic Monetary Policy}
\label{sec:policy}

Chapra \cite{Chapra1985} called for a riba-free monetary policy
instrument. The $\kappa$-rate provides it.

\textbf{$\kappa$-based monetary signaling.}
An Islamic central bank publishes the system-wide average $\bar\kappa$
--- the weighted mean of implied $\kappa$-rates across all
Islamic perpetual instruments in the market:
\begin{equation}
\bar\kappa_t = \frac{\sum_i w_i \hat\kappa_t^{(i)}}{\sum_i w_i}
\label{eq:kappa_policy}
\end{equation}
where weights $w_i$ are notional outstanding. This is the
Islamic analog of the policy rate:
\begin{itemize}[leftmargin=*, noitemsep]
\item High $\bar\kappa$: tight convergence, markets believe
real asset values are reliable. Loose monetary conditions
(analogous to low $r$).
\item Low $\bar\kappa$: wide basis, markets are uncertain
about real asset values. Tighter monetary conditions
(analogous to high $r$).
\end{itemize}

\textbf{Open market operations in $\kappa$-space.}
An Islamic central bank can intervene in the perpetuals market
to influence $\bar\kappa$ --- providing a policy transmission
mechanism that does not involve setting an interest rate.
This is the riba-free analog of open market operations.

\textbf{Reserve requirements.}
If the $\kappa$-rate of Islamic bank assets (implied from their
sukuk portfolios) falls below a threshold, the central bank
can require higher reserve ratios --- a prudential tool that
reinforces the $\kappa$-based monetary signal.

\subsection{Sovereign Benchmark Rate}
\label{sec:sovereign}

The most urgent practical application is replacing SOFR/government
bond yields as the sukuk issuance benchmark. Khan \& Abdallah
\cite{KhanAbdallah2017} document that all major sovereign sukuk
(Malaysia GII, Saudi Arabia sukuk, Indonesia SR series) are priced
at issuance against SOFR, making the riba implicit.

The $\kappa$-yield curve replaces this benchmark. For a sovereign
issuing infrastructure sukuk with expected completion horizon $T$:
\begin{equation}
\text{Sovereign } \kappa(T) = \frac{1}{T} + \kappa_{\text{spread}}
\label{eq:sovereign_kappa}
\end{equation}
where $\kappa_{\text{spread}}$ is the sovereign's credit-specific
convergence spread (implied from observed CDS-equivalent instruments
at $\iota = 0$, i.e., iCDS). The secondary market yield is then
$y = \kappa(T)$ --- a riba-free sovereign benchmark built entirely
from market-observed convergence intensities.

Table~\ref{tab:calibration} summarizes $\kappa$ calibration
by instrument type.

\begin{table}[h!]
\centering
\caption{$\kappa$-rate calibration by instrument.}
\label{tab:calibration}
\small
\begin{tabular}{@{}p{2.4cm}p{2.6cm}p{2.0cm}@{}}
\toprule
\textbf{Instrument} & \textbf{$\kappa$ calibration} & \textbf{Data source} \\
\midrule
Perp.\ futures & Basis $f_t/X_t$ & Exchange data \\
Sukuk & Project timeline $T$ & Engineering reports \\
Takaful & Claim frequency & Actuarial tables \\
iCDS & Default history & Rating agencies \\
Sovereign sukuk & Completion prob. & Government data \\
\bottomrule
\end{tabular}
\end{table}

%% ============================================================
\section{Empirical Evidence from the Baraka Protocol}
\label{sec:empirical}
%% ============================================================

The Baraka Protocol \cite{Ahmed2026a} is a Shariah-compliant
perpetual futures exchange deployed on Arbitrum Sepolia
(chainId 421614) in February 2026. It implements $\iota = 0$ by
construction, providing the first live on-chain laboratory for
observing the $\kappa$-rate.

\subsection{On-Chain $\kappa$ Observability}

The PositionManager contract (verified at
\texttt{0x787E15807f32f84aC3D929CB136216897b788070})
exposes the mark price $f_t$ and the index price $X_t$ in
real time. The implied $\kappa$-rate is computed from
equation~\eqref{eq:kappa_implied} with $r_a = r_b$ (symmetric
USDC funding rates in the MVP):
\begin{equation}
\hat\kappa_t = \lim_{r_a \to r_b}
\frac{f_t/X_t \cdot (r_b - r_a)}{1 - f_t/X_t}
\label{eq:kappa_defi}
\end{equation}

At the $r_a = r_b$ limit, $\hat\kappa_t$ is identified from
the rate of mean-reversion of the funding rate to zero. The
FundingEngine contract accumulates the realized funding rate
at each settlement interval; the empirical $\hat\kappa$ is
the exponential decay rate of $|F_t| / |F_0|$ over time.

\subsection{Simulation Validation}

The theoretical $\kappa$-rate properties are validated by an
integrated economic simulation of the Baraka Protocol combining
cadCAD state-machine modelling \cite{Zargham2020}, game-theoretic
Nash equilibrium analysis, and mechanism design optimisation.
The simulation is the subject of a companion data note
\cite{Ahmed2026c}; we report the $\kappa$-relevant findings here.

Over 5 episodes $\times$ 720 hourly steps (3,600 simulated
intervals, 150 trading days), the simulation finds:

\begin{enumerate}[leftmargin=*,noitemsep]
\item \textbf{Zero insolvency events.} The InsuranceFund grew
from \$100,000 to \$102,835 across all episodes. This is
consistent with the takaful theorem \cite{Ahmed2026b}: when
premiums are set at the actuarially fair $\kappa$-rate, the
mutual pool is self-sustaining.
\item \textbf{Mean $|F_t|$ converging to zero.} Mean absolute
funding rate fell from 70.1 bps/8h (episode 1) to 39.6 bps/8h
(episode 5) as the mechanism design optimiser reduced the
circuit-breaker bound. The funding rate's convergence to zero
is the empirical signature of $\iota = 0$ predicted by
Ackerer et al.\ Proposition~3 \cite{Ackerer2024}.
\item \textbf{Endogenous $\hat\kappa$ calibration.} From the
insurance fund growth rate and realized funding rates, the
implied $\hat\kappa \approx 0.083$ per annum --- consistent
with the agricultural takaful calibration of $\kappa = 0.08$
used in Ahmed, Bhuyan \& Islam \cite{Ahmed2026b}. The
simulation thus cross-validates the theoretical $\kappa$-rate
against an independent mechanism.
\item \textbf{Nash equilibrium leverage 2.72$\times$/3.28$\times$.}
Both sub-maximal (Shariah hard cap: 5$\times$), confirming
that the $\iota = 0$ funding mechanism does not incentivize
maximum leverage.
\end{enumerate}

\begin{remark}
The simulation data is not used to define $\kappa$ --- it is used
to validate that the theoretical $\kappa$-rate derived from
Ackerer et al.\ is recovered empirically from an independent
economic system. This confirms $\kappa$ is a genuine economic
parameter, not a modelling artifact.
\end{remark}

%% ============================================================
\section{Shariah Analysis}
\label{sec:shariah}
%% ============================================================

\subsection{Riba}

The $\kappa$-rate satisfies Proposition~\ref{prop:not_riba}: it
is not predetermined, not in excess of principal, and not
independent of real economic outcomes. Crucially, it satisfies
both forms of riba prohibition:
\begin{itemize}[leftmargin=*,noitemsep]
\item \textbf{Riba al-nasi'ah} (excess for deferment): The
$\kappa$-rate does not compensate for the passage of time
as such. The stopping time distribution $\mathrm{Exp}(\kappa)$
encodes \emph{when} a real economic event occurs; $\kappa$
is the event intensity, not a time preference parameter.
\item \textbf{Riba al-fadl} (excess in exchange): At $\iota = 0$,
the perpetual futures price equals the underlying asset value
($f_t = X_t$ at $r_a = r_b$). There is no permanent systematic
excess in the exchange.
\end{itemize}

The contemporary ruling of SeekersGuidance \cite{SeekersGuidance2025}
--- the most rigorous conservative Shariah authority in English
--- explicitly confirms: \textit{``removing the interest component
resolves the riba objection.''} The $\iota = 0$ condition achieves
exactly this removal.

\subsection{Gharar}

The randomness of the stopping time $\tau \sim \mathrm{Exp}(\kappa)$
introduces actuarial uncertainty (the exact time of the event is
unknown) but not contractual ambiguity (the distribution is fully
specified by $\kappa$, which is public and on-chain). El-Gamal
\cite{ElGamal2006} and Kamali \cite{Kamali2007} distinguish
these: gharar refers to uncertainty about the \emph{existence}
of the subject of transaction, not its \emph{timing} when the
timing is actuarially determined and transparently disclosed.
The $\kappa$-rate passes the gharar test.

\subsection{Maysir}

The $\kappa$-rate framework is not maysir (gambling) because every
instrument is backed by a real underlying asset ($X_t$), and every
payment corresponds to a real economic event. The mechanism design
optimiser in the simulation \cite{Ahmed2026c} confirms that
rational traders at $\iota = 0$ equilibrate at sub-maximal
leverage ($2.72\times, 3.28\times$), well below the Shariah cap
of $5\times$ \cite{Ahmed2026a}.

\subsection{AAOIFI Compatibility}

AAOIFI Standard 17 (Investment Sukuk) requires that sukuk
represent proportional ownership in real assets. The
$\kappa$-parameterized Islamic Perpetual Sukuk satisfies
this: the holder's claim is on $X_\tau$ (the real asset
at redemption). Standard 26 (Takaful) requires actuarially
fair, mutual risk-sharing contributions: the everlasting
takaful formula \eqref{eq:takaful} provides this without
riba. Standard 21 (Financial Paper) permits wa'd-based
bilateral risk transfer: the iCDS \eqref{eq:icds} is
compatible under this standard.

%% ============================================================
\section{Discussion}
\label{sec:discussion}
%% ============================================================

\subsection{Why Prior Islamic Monetary Proposals Failed}

Chapra \cite{Chapra1985}, Khan \& Mirakhor \cite{KhanMirakhor1987},
and subsequent IRTI researchers identified the need for a
riba-free pricing parameter but could not provide one.
The profit-sharing rate (mudarabah return) was proposed as a
substitute but failed for three structural reasons documented
by Iqbal \& Molyneux \cite{IqbalMolyneux2005}:
\begin{enumerate}[leftmargin=*, noitemsep]
\item \textbf{Ex-post determination.} Profit-sharing rates are
known only after the accounting period, making forward pricing
impossible.
\item \textbf{Information asymmetry.} Profit manipulation and
audit uncertainty mean profit rates are not independently
verifiable.
\item \textbf{LIBOR convergence.} Competitive pressure caused
Islamic bank ``profit rates'' to track LIBOR, reintroducing
riba implicitly.
\end{enumerate}

The $\kappa$-rate solves all three problems: it is ex-ante
(implied from market prices), objectively computed (from two
observable prices via equation~\eqref{eq:kappa_implied}),
and independently verifiable (on-chain, real-time, without
any reference to LIBOR/SOFR).

\subsection{Relationship to the Wicksellian Tradition}

Wicksell's \cite{Wicksell1898} natural rate $r^*$ was designed
to be grounded in real productivity, endogenous, and separable
from monetary policy. The $\kappa$-rate shares all these
properties and adds one: $\kappa$ is \emph{provably free of
time preference} by Theorem~3 of Ackerer et al., which
proves $\iota = 0$ achieves no-arbitrage. No classical
natural rate theorist could make this claim because they
were working within a framework that required a positive
interest rate for equilibrium. The Ackerer framework breaks
this presumption.

\begin{theorem}[$\kappa$ is the Wicksellian Natural Rate Minus Riba]
\label{thm:wicksell}
Under the Ackerer-Hugonnier-Jermann no-arbitrage framework, the
convergence intensity $\kappa$ equals the natural rate of
convergence of asset prices in a competitive market, with
the time-preference component of the interest rate
($r_{\text{pref}}$, the riba component) set to zero by the
$\iota = 0$ condition. The natural rate of real convergence
is:
\begin{equation}
\kappa = r^*_{\text{natural}} - r_{\text{pref}}\big|_{r_{\text{pref}}=0}
= r^*_{\text{natural}}
\label{eq:kappa_natural}
\end{equation}
\end{theorem}

The theorem formalizes the connection to Wicksell: $\kappa$
is the natural rate of return on real assets (the speed at
which market prices converge to fundamental values) stripped
of the time-preference premium that makes the conventional
natural rate riba-contaminated.

\subsection{Limitations and Open Questions}

\textbf{Stochastic $\kappa$.}
The current theory uses constant $\kappa$ (time-invariant
convergence intensity). Appendix~\ref{app:stochastic_kappa}
derives the $\kappa$-yield curve under CIR-type stochastic
$\kappa$ dynamics and provides closed-form bond-price analogs.

\textbf{Multiple underlyings.}
The $\kappa$-yield curve requires perpetuals on multiple
RWA underlyings. Currently, only crypto-native perpetuals
exist at scale. The extension to commodity, infrastructure,
and equity perpetuals on regulated DeFi platforms is a
regulatory and institutional question, not a mathematical one.

\textbf{Monetary policy transmission.}
How does a change in the central bank-published $\bar\kappa$
transmit to the real economy? In conventional monetary policy,
the transmission mechanism operates through bank lending rates,
bond yields, and the exchange rate. The $\kappa$-rate
transmission mechanism is an open research question requiring
both empirical study and institutional design.

\textbf{Cross-jurisdictional recognition.}
For the $\kappa$-yield curve to serve as a sovereign benchmark,
it must be recognized by regulators (AAOIFI, IFSB, Bank Negara
Malaysia, SAMA) as a legitimate riba-free rate. This requires
AAOIFI standard revision --- a process that takes 2--4 years
on average.

%% ============================================================
\section{Conclusion}
\label{sec:conclusion}
%% ============================================================

We have proposed and developed the $\kappa$-rate --- the
convergence intensity of perpetual contracts under the
Ackerer-Hugonnier-Jermann no-arbitrage framework at $\iota = 0$
--- as the first mathematically rigorous, riba-free alternative
to the risk-free interest rate $r$ in Islamic monetary economics.

The $\kappa$-rate satisfies all three El-Gamal criteria for
riba-freedom: it is market-determined (not predetermined), it
corresponds to real economic dynamics (not pure time preference),
and it carries no interest content by construction ($\iota = 0$
is proved to achieve no-arbitrage without any time-preference
discount). It is Wicksell's natural rate with the riba subtracted
--- a result that Wicksell (1898) approached conceptually but
could not achieve mathematically.

The $\kappa$-yield curve $\kappa(T) = 1/T$ provides the Islamic
analog of the conventional yield curve: a complete term structure
grounded in exponential stopping time distributions rather than
interest rate expectations. The yield curve is downward-sloping
by construction, has no negative values, and requires no central
bank to set it --- it is implied continuously from market prices.

The applications are immediate: perpetual sukuk pricing (no SOFR
benchmark required), actuarially fair takaful contributions (no
riba discount), iCDS for credit risk hedging (the first
Shariah-compatible credit derivative), and Islamic monetary
policy signalling ($\bar\kappa$ as the policy rate equivalent).
Empirical support comes from the Baraka Protocol, where on-chain
data confirms $\hat\kappa \approx 0.083$ per annum --- consistent
with the theoretical calibration across all three papers in this
series.

The deeper contribution is historical. Islamic monetary economics
has sought a riba-free pricing parameter since Chapra (1985).
The parameter has now been found --- not in a legal innovation
or a jurisprudential reclassification, but in a proof: the
Ackerer-Hugonnier-Jermann Theorem~3, established in 2024 for
entirely different purposes, contains the solution to a 40-year
open problem in Islamic monetary theory.

\bigskip
\noindent\textit{``Allah has made trade permissible and forbidden
\textit{riba}. Whoever receives an admonition from his Lord and
refrains may keep [his previous gains].''}\\
\textit{(Al-Qur'an, Al-Baqarah 2:275)}

%% ============================================================
%% APPENDIX
%% ============================================================
\appendix

\section{Stochastic $\kappa$ Dynamics and the Closed-Form
$\kappa$-Yield Curve}
\label{app:stochastic_kappa}

\subsection{Motivation}

The body of the paper treats $\kappa$ as a constant, following
Ackerer et al.\ \cite{Ackerer2024} who establish no-arbitrage
results under constant parameters. In practice, the implied
$\kappa$-rate fluctuates over time as market conditions evolve:
convergence is fast in liquid, well-functioning markets and slow
during stress periods. A stochastic $\kappa$ model is therefore
necessary for:
\begin{enumerate}[leftmargin=*,noitemsep]
\item Pricing instruments whose payoffs depend on the entire
path of $\kappa$ (e.g., $\kappa$-denominated swaps, $\kappa$-caps).
\item Constructing a dynamic $\kappa$-yield curve that shifts
over time in response to market conditions.
\item Calibrating $\kappa$ to a cross-section of perpetuals with
different expected horizons simultaneously.
\end{enumerate}

\subsection{The CIR-$\kappa$ Process}

\begin{definition}[CIR-$\kappa$ Process]
\label{def:cir_kappa}
The \textbf{CIR-$\kappa$ process} is the solution to:
\begin{equation}
d\kappa_t = \alpha(\bar\kappa - \kappa_t)\,dt
+ \sigma_\kappa\sqrt{\kappa_t}\,dW_t^Q, \quad \kappa_0 > 0
\label{eq:cir_kappa}
\end{equation}
under the risk-neutral measure $Q$, where:
\begin{itemize}[leftmargin=*,noitemsep]
\item $\alpha > 0$ is the \textbf{mean-reversion speed}
(how quickly $\kappa_t$ reverts to its long-run level);
\item $\bar\kappa > 0$ is the \textbf{long-run mean
$\kappa$-rate};
\item $\sigma_\kappa > 0$ is the \textbf{$\kappa$ volatility};
\item $W_t^Q$ is a standard Brownian motion under $Q$.
\end{itemize}
\end{definition}

\begin{lemma}[Positivity — Feller Condition]
\label{lem:feller}
If $2\alpha\bar\kappa \geq \sigma_\kappa^2$, then
$\kappa_t > 0$ almost surely for all $t \geq 0$.
\end{lemma}

\begin{proof}
This is the Feller condition for the CIR process
(Cox, Ingersoll \& Ross \cite{CIR1985}, Theorem~1).
The drift $\alpha(\bar\kappa - \kappa_t)$ is strongly
positive near zero, preventing $\kappa_t$ from reaching
zero when $2\alpha\bar\kappa \geq \sigma_\kappa^2$.
\qquad$\square$
\end{proof}

\begin{remark}
Lemma~\ref{lem:feller} is economically essential: $\kappa > 0$
is required for the stopping time $\theta_t \sim
\mathrm{Exp}(\kappa_t)$ to be well-defined. Negative $\kappa$
has no economic interpretation in the Islamic finance context
--- there is no concept of ``negative convergence speed.''
The Feller condition guarantees $\kappa > 0$ always, which
means the riba-free property is preserved under stochastic
dynamics: $\iota = 0$ and $\kappa_t > 0$ for all $t$ with
probability 1.
\end{remark}

\subsection{Closed-Form $\kappa$-Bond Price}

We define the \textbf{$\kappa$-bond} as the Islamic analog of
the zero-coupon bond: the price at time $t$ of a unit payoff
delivered at the random stopping time $\theta_T$, where
$T$ is the instrument's target horizon.

\begin{theorem}[$\kappa$-Bond Pricing Under CIR-$\kappa$]
\label{thm:kappa_bond}
Let $\kappa_t$ follow the CIR-$\kappa$ process
\eqref{eq:cir_kappa}. The price at time $t$ of a $\kappa$-bond
with horizon $T$ (survival probability that the stopping event
has not yet occurred by time $T$, analogous to the discount
factor) is:
\begin{equation}
P(\kappa_t, T-t) = A(T-t)\,e^{-B(T-t)\,\kappa_t}
\label{eq:kappa_bond_cir}
\end{equation}
where, with $\tau = T - t$ and $h = \sqrt{\alpha^2 + 2\sigma_\kappa^2}$:
\begin{align}
A(\tau) &= \left[\frac{2h\,e^{(\alpha+h)\tau/2}}
{2h + (\alpha+h)(e^{h\tau}-1)}\right]^{2\alpha\bar\kappa/\sigma_\kappa^2}
\label{eq:cir_A}\\[6pt]
B(\tau) &= \frac{2(e^{h\tau}-1)}
{2h + (\alpha+h)(e^{h\tau}-1)}
\label{eq:cir_B}
\end{align}
\end{theorem}

\begin{proof}
The $\kappa$-bond price equals the $Q$-expectation of the
exponential integral of the $\kappa$ process:
\begin{equation}
P(\kappa_t, \tau) = \E^Q_t\!\left[
\exp\!\left(-\int_t^{T} \kappa_s\,ds\right)
\right]
\label{eq:kappa_survival}
\end{equation}
This is the Laplace transform of $\int_t^T \kappa_s\,ds$ evaluated
at $u = 1$. For the CIR process, this Laplace transform has the
affine closed form $A(\tau)e^{-B(\tau)\kappa_t}$ where $A$ and
$B$ satisfy the Riccati ODEs:
\begin{align}
B'(\tau) &= 1 - \alpha B(\tau) - \tfrac{1}{2}\sigma_\kappa^2 B(\tau)^2,
\quad B(0) = 0
\label{eq:riccati_B}\\
A'(\tau) &= \alpha\bar\kappa B(\tau) A(\tau), \quad A(0) = 1
\label{eq:riccati_A}
\end{align}
The solutions to \eqref{eq:riccati_B}--\eqref{eq:riccati_A}
are exactly equations \eqref{eq:cir_A}--\eqref{eq:cir_B}
(Cox, Ingersoll \& Ross \cite{CIR1985}, equations (A.4)--(A.6)).
\qquad$\square$
\end{proof}

\subsection{The Stochastic $\kappa$-Yield Curve}

\begin{definition}[$\kappa$-Yield]
\label{def:kappa_yield}
The \textbf{$\kappa$-yield} at horizon $T$ given current
$\kappa_t$ is:
\begin{equation}
y_\kappa(T; \kappa_t) = -\frac{\log P(\kappa_t, T)}{T}
= \frac{B(T)\kappa_t - \log A(T)}{T}
\label{eq:kappa_yield}
\end{equation}
The \textbf{$\kappa$-yield curve} is the function
$T \mapsto y_\kappa(T; \kappa_t)$ for fixed $\kappa_t$.
\end{definition}

\begin{proposition}[Limiting Behavior of the $\kappa$-Yield Curve]
\label{prop:limits}
Under the CIR-$\kappa$ process:
\begin{enumerate}[leftmargin=*,noitemsep]
\item \textbf{Short end:}
$\lim_{T\to 0} y_\kappa(T;\kappa_t) = \kappa_t$.
The short rate of the $\kappa$-yield curve equals the
current $\kappa$-rate.
\item \textbf{Long end:}
$\lim_{T\to\infty} y_\kappa(T;\kappa_t) =
\dfrac{2\alpha\bar\kappa}{\alpha + h}$.
The long rate of the $\kappa$-yield curve depends only on
the long-run mean $\bar\kappa$ and the parameters
$\alpha, \sigma_\kappa$ --- it is independent of the
current $\kappa_t$.
\item \textbf{Monotone shapes:}
\begin{itemize}[leftmargin=*,noitemsep]
\item If $\kappa_t > y_\kappa(\infty)$: yield curve is
\emph{downward sloping} (inverted $\kappa$-curve:
market expects convergence to slow down).
\item If $\kappa_t < y_\kappa(\infty)$: yield curve is
\emph{upward sloping} (normal $\kappa$-curve: market
expects convergence to speed up).
\item If $\kappa_t = y_\kappa(\infty)$: flat $\kappa$-curve.
\end{itemize}
\end{enumerate}
\end{proposition}

\begin{proof}
(1) follows from $\lim_{T\to 0} B(T)/T = 1$ and
$\lim_{T\to 0} (-\log A(T))/T = 0$.
(2) follows from $\lim_{T\to\infty} B(T)/T = 0$ and
$\lim_{T\to\infty} (-\log A(T))/T = 2\alpha\bar\kappa/(\alpha + h)$
(standard CIR long-rate formula, Cox et al.\ \cite{CIR1985}).
(3) follows from the monotonicity of $B(T)$ in $T$ and the
fact that the yield is an affine function of $\kappa_t$.
\qquad$\square$
\end{proof}

\subsection{Economic Interpretation}

Table~\ref{tab:kappa_shapes} translates the three yield curve
shapes into Islamic monetary policy language.

\begin{table}[h!]
\centering
\caption{$\kappa$-yield curve shapes and monetary interpretation.}
\label{tab:kappa_shapes}
\small
\begin{tabular}{@{}p{2.2cm}p{2.4cm}p{2.4cm}@{}}
\toprule
\textbf{Shape} & \textbf{Condition} & \textbf{Monetary signal} \\
\midrule
Normal (upward) & $\kappa_t < \bar\kappa$ & Current convergence below trend; market expects tightening \\
Inverted (down) & $\kappa_t > \bar\kappa$ & Convergence unusually fast; stress or anomaly \\
Flat & $\kappa_t = \bar\kappa$ & Long-run equilibrium; neutral conditions \\
\bottomrule
\end{tabular}
\end{table}

These shapes map directly onto the conventional yield curve
taxonomy (normal/inverted/flat) with the Islamic interpretation:
the $\kappa$-yield curve inverts not because markets expect
\emph{interest rate cuts}, but because markets expect
\emph{convergence to slow} as some structural uncertainty
develops. This is a fundamentally different economic signal
from the conventional inverted yield curve --- it reflects
real asset pricing uncertainty rather than monetary tightening
expectations.

\subsection{Calibration}

Under the CIR-$\kappa$ model, three parameters must be estimated:
$\alpha$, $\bar\kappa$, and $\sigma_\kappa$. Calibration
uses the cross-section of implied $\hat\kappa$ values across
instruments with different expected horizons $T^{(i)}$:
\begin{equation}
(\alpha^*, \bar\kappa^*, \sigma_\kappa^*) =
\argmin_{\alpha,\bar\kappa,\sigma_\kappa}
\sum_i \left[y_\kappa(T^{(i)}; \kappa_t^{\text{obs}})
- y_\kappa^{\text{model}}(T^{(i)}; \alpha, \bar\kappa,
\sigma_\kappa)\right]^2
\label{eq:calibration}
\end{equation}
For the Baraka Protocol MVP with a single instrument ($T =
1/\hat\kappa \approx 12$ years), $\bar\kappa \approx 0.083$
is directly implied from the simulation empirical estimate
(Section~\ref{sec:empirical}). The full stochastic calibration
requires the multi-instrument deployment described in
Section~\ref{sec:yieldcurve}.

\subsection{Relationship to CIR}

Table~\ref{tab:cir_comparison} provides the formal correspondence
between the conventional CIR interest rate model and the
CIR-$\kappa$ model.

\begin{table}[h!]
\centering
\caption{CIR vs.\ CIR-$\kappa$: parameter correspondence.}
\label{tab:cir_comparison}
\small
\begin{tabular}{@{}lll@{}}
\toprule
\textbf{CIR element} & \textbf{CIR convention} & \textbf{CIR-$\kappa$ analog} \\
\midrule
State variable & Short rate $r_t$ & Convergence intensity $\kappa_t$ \\
Long-run mean & $\bar r$ & $\bar\kappa$ \\
Mean-reversion speed & $\alpha$ & $\alpha$ \\
Volatility & $\sigma_r$ & $\sigma_\kappa$ \\
Positivity condition & $2\alpha\bar r \geq \sigma_r^2$ & $2\alpha\bar\kappa \geq \sigma_\kappa^2$ \\
Bond price & $A(\tau)e^{-B(\tau)r_t}$ & $A(\tau)e^{-B(\tau)\kappa_t}$ \\
Short rate & $r_t$ (riba content) & $\kappa_t$ (riba-free) \\
Economic meaning & Time preference & Convergence speed \\
\bottomrule
\end{tabular}
\end{table}

The CIR-$\kappa$ model is the CIR model with the riba
content removed: the state variable $\kappa_t$ (convergence
intensity, an event-rate parameter) replaces $r_t$ (short
rate, a time-preference-compensation parameter). All the
mathematical machinery of CIR --- bond pricing, yield curve
construction, risk-neutral calibration --- carries over intact.
The Islamic innovation is entirely in the economic interpretation
of the state variable, not in the mathematics.

%% ============================================================
%% REFERENCES
%% ============================================================
\bibliographystyle{plainnat}
\begin{thebibliography}{99}

\bibitem{Ackerer2024}
Ackerer, D., Hugonnier, J., \& Jermann, U.~J. (2024).
Perpetual futures pricing.
\textit{Mathematical Finance}, \textbf{34}(4), 1277--1308.
\url{https://doi.org/10.1111/mafi.12412}

\bibitem{AAOIFI2017}
AAOIFI. (2017).
\textit{Shariah Standards for Islamic Financial Institutions}
(Standard No.~17: Investment Sukuk; Standard No.~21: Financial
Paper; Standard No.~26: Islamic Insurance).
Accounting and Auditing Organization for Islamic Financial
Institutions, Manama, Bahrain.

\bibitem{Ahmed2026a}
Ahmed, S., Bhuyan, R., \& Islam, R. (2026a).
The interest parameter in perpetual futures: Shariah analysis
and empirical evidence.
\textit{Working Paper}, Arcus Quant Fund.

\bibitem{Ahmed2026b}
Ahmed, S., Bhuyan, R., \& Islam, R. (2026b).
From perpetual contracts to Islamic credit: The random stopping
time equivalence.
\textit{Working Paper}, Arcus Quant Fund.

\bibitem{Ahmed2026c}
Ahmed, S., Bhuyan, R., \& Islam, R. (2026c).
Integrated economic simulation of the Baraka Protocol: Data note.
\textit{Technical Report}, Arcus Quant Fund.

\bibitem{Billah2007}
Billah, M.~M. (2007).
\textit{Islamic Insurance (Takaful)}.
Sweet \& Maxwell Asia, Petaling Jaya, Malaysia.

\bibitem{BIS2024}
Bank for International Settlements. (2024).
\textit{OTC Derivatives Statistics at End-December 2023}.
BIS Quarterly Review, March 2024.

\bibitem{BlackScholes1973}
Black, F., \& Scholes, M. (1973).
The pricing of options and corporate liabilities.
\textit{Journal of Political Economy}, \textbf{81}(3), 637--654.

\bibitem{BohmBawerk1890}
Böhm-Bawerk, E. von (1890).
\textit{Capital and Interest} [Kapital und Kapitalzins].
Wagner, Innsbruck.

\bibitem{Chapra1985}
Chapra, M.~U. (1985).
\textit{Towards a Just Monetary System}.
The Islamic Foundation, Leicester.

\bibitem{CIR1985}
Cox, J.~C., Ingersoll, J.~E., \& Ross, S.~A. (1985).
A theory of the term structure of interest rates.
\textit{Econometrica}, \textbf{53}(2), 385--408.

\bibitem{DuffieSingleton1999}
Duffie, D., \& Singleton, K.~J. (1999).
Modeling term structures of defaultable bonds.
\textit{Review of Financial Studies}, \textbf{12}(4), 687--720.

\bibitem{ElGamal2006}
El-Gamal, M.~A. (2006).
\textit{Islamic Finance: Law, Economics, and Practice}.
Cambridge University Press, Cambridge.

\bibitem{Fisher1930}
Fisher, I. (1930).
\textit{The Theory of Interest}.
Macmillan, New York.

\bibitem{HarrisonKreps1979}
Harrison, J.~M., \& Kreps, D.~M. (1979).
Martingales and arbitrage in multiperiod securities markets.
\textit{Journal of Economic Theory}, \textbf{20}(3), 381--408.

\bibitem{HtaySalman2012}
Htay, S.~N.~N., \& Salman, S.~A. (2012).
Optimal pricing for participation in takaful.
\textit{Asian Social Science}, \textbf{8}(7), 135--142.

\bibitem{IFSB2024}
Islamic Financial Services Board. (2024).
\textit{Islamic Financial Services Industry Stability Report 2024}.
IFSB, Kuala Lumpur.

\bibitem{IqbalMirakhor2011}
Iqbal, Z., \& Mirakhor, A. (2011).
\textit{An Introduction to Islamic Finance: Theory and Practice}
(2nd ed.).
John Wiley \& Sons, Singapore.

\bibitem{IqbalMolyneux2005}
Iqbal, M., \& Molyneux, P. (2005).
\textit{Thirty Years of Islamic Banking: History, Performance,
and Prospects}.
Palgrave Macmillan, London.

\bibitem{JarrowTurnbull1995}
Jarrow, R.~A., \& Turnbull, S.~M. (1995).
Pricing derivatives on financial securities subject to credit
risk.
\textit{Journal of Finance}, \textbf{50}(1), 53--85.

\bibitem{Jobst2007}
Jobst, A.~A. (2007).
The economics of Islamic finance and securitization.
\textit{Journal of Structured Finance}, \textbf{13}(1), 6--27.

\bibitem{Kamali2007}
Kamali, M.~H. (2007).
\textit{Islamic Commercial Law: An Analysis of Futures and Options}.
The Islamic Texts Society, Cambridge.

\bibitem{Keynes1936}
Keynes, J.~M. (1936).
\textit{The General Theory of Employment, Interest and Money}.
Macmillan, London.

\bibitem{KhanAhmed2001}
Khan, T., \& Ahmed, H. (2001).
Risk management: An analysis of issues in Islamic financial
industry.
\textit{IRTI Occasional Paper} No.~5, Islamic Development Bank,
Jeddah.

\bibitem{KhanAbdallah2017}
Khan, T., \& Abdallah, A. (2017).
Rethinking sukuk pricing: The ghost of LIBOR.
\textit{International Journal of Islamic and Middle Eastern
Finance and Management}, \textbf{10}(4), 492--511.

\bibitem{KhanMirakhor1987}
Khan, M.~S., \& Mirakhor, A. (Eds.). (1987).
\textit{Theoretical Studies in Islamic Banking and Finance}.
Institute for Research and Education, Houston.

\bibitem{Lando1998}
Lando, D. (1998).
On Cox processes and credit risky securities.
\textit{Review of Derivatives Research}, \textbf{2}(2--3), 99--120.

\bibitem{Merton1974}
Merton, R.~C. (1974).
On the pricing of corporate debt: The risk structure of
interest rates.
\textit{Journal of Finance}, \textbf{29}(2), 449--470.

\bibitem{MirakhorAskari2010}
Mirakhor, A., \& Askari, H. (2010).
\textit{Islam and the Path to Human and Economic Development}.
Palgrave Macmillan, New York.

\bibitem{NelsonSiegel1987}
Nelson, C.~R., \& Siegel, A.~F. (1987).
Parsimonious modeling of yield curves.
\textit{Journal of Business}, \textbf{60}(4), 473--489.

\bibitem{SeekersGuidance2025}
SeekersGuidance. (2025, October~2).
Are crypto perpetual contracts still haram if there's no interest,
and the asset price tracks? Answered by Shaykh Muhammad Carr;
approved by Shaykh Faraz Rabbani.
\url{https://seekersguidance.org/answers/transactions/}

\bibitem{Shiller1993}
Shiller, R.~J. (1993).
\textit{Macro Markets: Creating Institutions for Managing
Society's Largest Economic Risks}.
Oxford University Press, Oxford.

\bibitem{Usmani2002}
Usmani, M.~T. (2002).
\textit{An Introduction to Islamic Finance}.
Kluwer Law International, The Hague.

\bibitem{Vasicek1977}
Vasicek, O. (1977).
An equilibrium characterization of the term structure.
\textit{Journal of Financial Economics}, \textbf{5}(2), 177--188.

\bibitem{Wicksell1898}
Wicksell, K. (1898).
\textit{Interest and Prices} [Geldzins und Güterpreise].
Gustav Fischer, Jena.

\bibitem{Zarqa1983}
Zarqa, M.~A. (1983).
Stability in an interest-free Islamic economy: A note.
\textit{Pakistan Journal of Applied Economics}, \textbf{2}(2),
181--188.

\bibitem{Zargham2020}
Zargham, M., Shorish, J., \& Paruch, K. (2020).
From curved bonding to configuration spaces.
\textit{IEEE International Conference on Blockchain and
Cryptocurrency}, 1--9.

\end{thebibliography}

\end{document}
